  \documentclass{article}
\usepackage{amsmath}
\usepackage{graphicx}
\graphicspath{ {images/} }
\usepackage[spanish]{babel}
\usepackage{bbm}
\usepackage{amsthm}
\usepackage{authblk}
\usepackage{hyperref}
\usepackage[utf8]{inputenc}
\usepackage{mathrsfs}
\usepackage{amsfonts}
\usepackage{amssymb}
\usepackage{enumerate}
\usepackage[left=3cm,right=2cm,top=2cm,bottom=2cm]{geometry}
\usepackage{epsfig}
 \renewcommand{\baselinestretch}{0.93}
 \thispagestyle{empty}
 \pagestyle{empty}
\setlength{\textheight}{53pc} \setlength{\textwidth} {36pc}
\setlength{\topmargin}{-4mm}
\usepackage{multicol}
\newcommand*{\email}[1]{%
    \normalsize\href{mailto:#1}{#1}\par
    }
\setlength{\columnsep}{1cm}
\newcommand{\N}{\mathbb{N}}
\newcommand{\calM}{\cal{M}}
\newcommand{\calP}{\cal{P}}
\newcommand{\F}{\mathbb{F}}
\newcommand{\R}{\mathbb{R}}
\newcommand{\Z}{\mathbb{Z}}
\newcommand{\Q}{\mathbb{Q}}
\newcommand{\C}{\mathbb{C}}
\newcommand{\bbH}{\mathbb{H}}


\theoremstyle{definition}
\newtheorem{ejer}{Ejercicio}
\author{Marcos Morales}
\affil{Universidad Finis Terrae}
\title{Primera Semana\\}



	


\begin{document}

\begin{center}
    \Large{\textbf{Taller con décimas}}
\end{center}


\setlength{\parindent}{0cm}


\textbf{Ejercicio 1:} Considere los planos $\displaystyle \Pi_1: 2x-y+z-3=0$ y  $\displaystyle \Pi_2: x+y-2z-4=0$. Determine la recta intersección $\Pi_1 \cap \Pi_2$ y escriba sus ecuaciones general y paramérica \\ \\ \\





\textbf{Ejercicio 2:} Calcule los siguientes límites, si no existen demuéstrelo \\

a) $\displaystyle \lim_{(x,y)\to (0,0)} \frac{x^2 - y^2}{x^2 + y^2}$ \\ \\

b) $\displaystyle \lim_{(x,y)\to (0,0)} \frac{1 - \cos(x^2+y^2)}{x^2 + y^2}$ \\ \\ \\





\textbf{Ejercicio 3:} Considere la superficie $f(x,y)= kx^3-2ky$. Determine el valor de $k$ para que el plano tangente a $f$ en el punto $(x,y)=(1,1)$, sea perpendicular al plano $3x-2y+z-4=0$



\end{document} 